\section{Relationship between \piI and \AOGS}
\label{s:encoGames}
             

To compare the \AOGS and \piI representations of the (call-by-value)
$\lambda$-calculus, we set a mapping from \AOGS configurations and environments 
to  \piI processes.
The mapping is reported in Figure~\ref{f:OGS_pi_cbv}. 
%It maps OGS configurations into \piI configurations (cf., Section~\ref{sss:os_piI}).
It is an extension of Milner's encoding of the $\lambda$-calculus and is
therefore indicated with the same symbol  $\cbvSymb$.
The optimisation $\qencV$ of Milner's encoding mentioned at the end of
the previous section is correspondingly extended to an optimisation of the mapping for
configurations. Again, we shall sometimes make use of such optimisations in proofs. 

The mapping uses a representation of environments $\gamma$ as associative lists, introduced in
Section~\ref{s:nota}.
So it relies on a choice of an order on the names that form the domain of $\gamma$.
This choice is irrelevant in the definition of the mapping into \piI, as behavioural equivalence in
\piI is invariant under commutativity and associativity of parallel composition.

\begin{remark}
 The  encoding of a configuration $\FF$ with name-support $\phi$ 
does not depend on $\phi$.
This name-support $\phi$ is used in OGS both to enforce freshness of names,
and to deduce the polarity of names,
as represented by the function $\Label$.
And indeed, the process $\encoConI{\FF}$ has its set of free names included
in $\phi$, and uses $P$-names in outputs and $O$-names in inputs.
The freshness property could be made more explicit by employing configurations also in \piI.
The polarity property could also be stated in $\pi$-calculus using 
i/o-sorting~\cite{PiSa96b}.
Indeed, a correspondence between arenas of game semantics (used to enforce polarities of moves) 
and sorting has  been 
explored~\cite{10.1145/224164.224189,10.1016/S0304-3975-99-00039-0}.
\end{remark}


\begin{figure}
%\Mybar 
Encoding of environments: \hfill   $ $ %  \dotfill $ $ 
\[ 
\begin{array}{rcl}
\encoEnvI{\pmap{y}{V}\cdot\gamma'  } & \defeq & \encoIVa {V} y | \encoEnvI{\gamma'  }
\\[\myptD] 
\encoEnvI{\pmap{q}{(E,p)}\cdot\gamma'  } & \defeq & 
\inp qx .\encoIa {E[x]} p  | \encoEnvI{\gamma'  } \\[\myptD]
\encoEnvI{ \emptymap} & \defeq & \nil
\end{array} 
\]
Encoding of configurations: \hfill  $ $ % \dotfill $ $ 
\[ 
\begin{array}{rcl}
%Encoding of active configurations $\conf{M,p,\gamma}$: 
\encoConI{ \conf{M,p,\gamma, \phi}}
& \defeq &   %\pairP { 
\encoIa {M} p | \encoEnvI \gamma 
%} 
\\[\myptD]
% Encoding of passive configurations $\conf{\gamma}$: 
 \encoConI{ \conf{\gamma, \phi}}& \defeq & %\pairP { 
 \encoEnvI \gamma  
% }
\\[\myptD]
% Encoding of the initial configuration:
 \encoConI{ \initconf{\phi}{M}} & \defeq &  %\pairP {
 \encoI   M
 %} 
\end{array}
 \] 
\caption{From OGS environments and configurations to \piI }
\label{f:OGS_pi_cbv}
%\Mybar 
\end{figure} 








\subsection{Operational correspondence}
\label{s:OC}

The following theorems establish the operational correspondence between the \AOGS and \piI representations. 

\begin{theorem} %[from OGS to \piI]
\label{t:opcorrCON_w}
$ $ \begin{enumerate}

\item If $\FF \LongrightarrowA \FF'$, then 
$\encoConI \FF \LongrightarrowPI  \; \contr \encoConI {\FF'}$;


\item If $\FF \ArrA\ell \FF'$, then
$ \encoConI \FF \ArrPI\ell  
\; \wbPI
\encoConI {\FF'}$.
\end{enumerate} 
\end{theorem} 



\begin{theorem} %[from \piI to  OGS]
\label{t:opcorrCON2_w}
$ $ \begin{enumerate}


\item If 
$ \encoConI \FF \LongrightarrowPI \PP$ then there is $\FF'$ such that 
$\FF \LongrightarrowA \FF'$  and  
$\PP  \wbPI \encoConI {\FF'}$; 


\item If
$ \encoConI \FF \ArrPI\ell \PP$
and  $\ell$ is an output, 
 then there is $\FF'$ such that
 $\FF \ArrA\ell \FF'$
and $\PP \wbPI  \encoConI {\FF'}$;


\item If  $\FF$ is passive and 
$ \encoConI \FF \ArrPI{\ell} \PP$, 
then 
 there is $\FF'$ such that 
$\FF \ArrA{\ell} \FF'$
and $\PP  \wbPI \encoConI {\FF'}$.
\end{enumerate} 
\end{theorem} 

\noindent
Both theorems above,  as well as similar operational correspondence theorems in this and next
section,  are first proved in terms of the optimised encoding  $\qencV$ 
(introduced in \cite{DurierHS18} and discussed in the previous section and in 
Appendix~\ref{a:oc}), 
and exploiting its
tight correspondence with the main encoding $\cbvSymb$. 
Moreover, 
  the operational correspondence 
is first established 
on strong
transitions,  along the lines of analogous correspondence results, such as  those detailed  in
Section~\ref{subsec:seq-ts-pi}  and Appendix~\ref{a:oc}. 
The correspondence on strong transitions, as usual in these cases, 
is carried out by  induction on the depth of the proof of a given transition; then the
 results are extended to weak transitions using induction on the length of the  given weak
 transition.  
 The advantage of   $\qencV$ is that it does not have initial
`administrative' reductions; hence any initial reduction corresponds to an internal move
of a $\lambda$-term.   This, for instance, allows us to have 
 the expansion relation  $\contr$ (in place of the coarser $\wbPI$, even though the final
 goal is  the result for  $\wbPI$), 
in the  statement about silent actions for the 
operational correspondence from \piI to \AOGS 
(see,  e.g., Lemma~\ref{l:opcorrCON2} in the Appendix, and similar lemmas). 
The use of expansion 
is
essential, both to derive the  operational correspondence theorems  on weak transitions,  and to use
the theorems in up-to techniques for  $\piI$ (more generally, in applications of the
theorems in which one reasons about the number of steps performed). 


\reviewComment{
Theorems 6.2 and 6.3: this seems to be at the core of the
translation. Could we have some details?}
\reviewCommentDS{We have added a paragraph following the theorems. The paper contains
  several operational correspondence results, whose proofs are not particularly
  interesting, and follow rather standard induction arguments.  These specific results are
similar to operational correspondence results that are found in papers that study the
relationship between lambda terms and their encoding pi-calculus  terms }


In Theorem~\ref{t:opcorrCON2_w}, a  clause is 
 missing
 for
input actions from $ \encoConI \FF$ when 
 $\FF$ active. Indeed such actions are possible in \piI, stemming from the (encoding of
 the) environment of $\FF$, whereas they are not possible in \AOGS. This is rectified in
 Section~\ref{subsec:seq-ts-pi}, introducing a constrained LTS for \piI, and in Section~\ref{s:cogs}, 
 considering a concurrent OGS.   
 
\begin{corollary}
\label{c:tracesGSpi}
If $\FF \ArrA s$ then also $ \enca \FF \ArrPI s $.
\end{corollary} 

\begin{proof}
By induction on the length of trace $s$.
\end{proof} 

\subsection{An \outputpr Transition System}
\label{subsec:seq-ts-pi}

We define an LTS for \piI in which input actions are
visible only if no output can be consumed, 
either as a  visible action or through an internal action 
(i.e., syntactically the process has no unguarded output). 
The new LTS, called \emph{\outputpr} and indicated as \opLTS,  
is defined on top of the ordinary one 
by means of the two rules below. A process $P$ is \emph{input reactive} if 
whenever $\PP \arrPI \mu  P'$, for some $\mu,P'$ 
 then  $\mu$ is an input action.
 \[
\begin{array}{l}
  \inferrule{\PP \arrPI{ \mu} \PP' \andalso  
  \mbox{$\PP$ input reactive}
 }{\PP \arrN {\mu} \PP'}
%\mbox{$\mu'$ is an input action}
\\
  \inferrule{\PP \arrPI{ \mu} \PP'
 }{\PP \arrN {\mu} \PP'}
\mbox{$\mu$ is an output or $\tau$  action}
\end{array}
 \]
% The LTS gives priority to output and silent actions, which forbid
% input actions. 
%We call 
The \opLTS captures an aspect of   \emph{sequentiality} 
in  $\pi$-calculi:  
a  free input prefix  is to be  thought of as a service offered to the external environment; 
in a sequential system, such a service is available only if 
there are no ongoing computations due to previous interrogations of the server. 
An ongoing computation is represented by  
a $\tau$-action, indicating a step  of computation  internal to the server, or 
an output, indicating either an answer to a client  or a request to an
external server.  
The constraint imposed by the 
new LTS can also be formalised compositionally,
as done in Figure~\ref{t:pii_op}.

\begin{figure*}[ht]
%\begin{table} 
\begin{center}
\begin{tabular}{rlcrl}
{\trans{alpha}}: &  $\displaystyle{  P\equival P' \andalso   { P'} \arrN\mu    { P''}
  \over 
  P \arrN\mu {P''}  
 } $ &  
% }
{\trans{ pre}}:&  
$\displaystyle{ %\bn \mu = \tilx \andalso \tilx \cap \phi = \emptyset  
\over
{\mu.      P} \arrN\mu {P}}$   
 \\[\mysp]
\multicolumn{5}{c}{ 
  {\trans{ parL}}:\; \; $\displaystyle{   P \arrN\mu { P'}
 \andalso  
  \mbox{($Q$ input reactive) or ($\mu$ is an output or  $\tau$)}
 \over  { P | Q}   \arrN\mu
{P'| Q} } $  }
\\[\mysp] 
\multicolumn{5}{c}{ 
    {\trans{com}}:  \; \;    $\displaystyle{ {P} \arrN{ \mu }{P'} \andalso   
 Q
\arrN{\outC\mu}{Q'}
\andalso \mu\neq \tau,   \; \tilx = \bn\mu 
  \over
{     P |  Q }\arrN{ \tau} 
{
 \res{  \tilx} (P'
|  Q') } }$    
}
 \\[\mysp]
{\trans{ res}}:& $\displaystyle{{  P} \arrN{\mu}{P'}
\andalso  x \not \in \n\mu
 \over
{ \res x     P}   \arrN{ \mu}{\res x P'  }} $ 
&
    {\trans{rep}}: &
 $\displaystyle{ {\mu.  P }\arrN{\mu} {P'}
 \over
{ !\mu. P }   \arrN{ \mu}{ P' | !\mu. P } } $  
 \\[\mysp]
\multicolumn{5}{c}{ 
    {\trans{con}}: \;\;
$  \displaystyle{   P \arrN \mu {P'}
\andalso  \mbox{ } {\rmmm K}  \defi
  \abs{\til y} Q  
\andalso 
\mbox{ } \abs{\til y} Q \equival \abs{\til x} P
    \over 
{
 \app {\rmmm K}{\til x}}  \arrN\mu 
{P'} }$
}\end{tabular} 
\end{center} 
\caption{A compositional presentation of the \outputpr LTS  for \piI}
\label{t:pii_op}
 \Mybar 
 %  \end{table}
\end{figure*}

\reviewComment{rule "pre": what is $\phi$?}
\reviewCommentDS{removed (probably a wrong cut-and-paste)}
Under the  \opLTS,  the analogue of Theorem~\ref{t:opcorrCON_w} continues to hold: 
in  \AOGS configurations,
input transitions only occur in passive configurations, 
and the encodings of passive configurations are input-reactive processes. 
However, now we have   the full  converse  of Theorem~\ref{t:opcorrCON2_w} and, as a
consequence, we can also establish  the converse direction of 
Corollary~\ref{c:tracesGSpi}.

\DSOCTb

\begin{lemma} %[from \piI to  OGS, sequential LTS]
\label{l:opcorrCON3_w}
% \begin{theorem} %[from \piI to  OGS, sequential LTS]
% \label{t:opcorrCON3_w}
$ $ \begin{enumerate}

% \item  If
% $\encoConV{\FF } \arrN  {\initName (p)} \PP$, then there is $\FF'$ 
% such that
%  $\FF \arr{\initName (p)} \FF'$ 
% and $\PP = \encoConV{\FF' }$ ;

\item If 
$ \encoConI \FF \ArrN {} \PP$ then there is $\FF'$ such that 
$\FF \LongrightarrowA \FF'$  and  
$\PP  \wbPI \encoConI {\FF'}$ ; 


\item If
$ \encoConI \FF \ArrN\ell \PP$
 then there is $\FF'$ such that
 $\FF \ArrA\ell \FF'$
and $\PP \wbPI \encoConI {\FF'}$.
\end{enumerate} 
%\end{theorem} 
\end{lemma} 

\noindent
Again, in the proof,  we first establish a correspondence result on strong transitions,
and reason on the optimised encoding $\qencV$. 
See Appendix~\ref{a:oc} for details.


Lemma~\ref{l:opcorrCON3_w}  
talks about the \opLTS of \piI;
  however the theorems make use of the
ordinary expansion relation  $\expa $, which is defined on the ordinary LTS. 
We now show that 
such uses of
expansion can actually be replaced by expansion on the \opLTS, defined as ordinary
expansion, but on the \opLTS.

A behavioural relation on the ordinary LTS `almost always' implies the corresponding
relation on the \opLTS.  The only exceptions are produced by processes that are related
although they exhibit different divergence behaviours (by divergence we mean the
possibility of performing an infinite number of $\tau$ actions) and at the same time may
perform input actions. For instance, using $\tau^\omega$ to indicate a purely divergent
process  (i.e., a process whose only transition is $\tau^\omega \arr\tau \tau^\omega$), 
 the processes $a.\nil | \tau^\omega $ and $a.\nil$ are  bisimilar (and trace equivalent)
 in the ordinary LTS, but they are not so in the \opLTS, since the divergence in the former
 process prevents observation of the input at $a$. 


Indeed, a `bisimilarity respecting divergence' would imply bisimilarity on the \opLTS. 
Such  `bisimilarity respecting divergence', written $\wbPIDIV$, is defined as $\wbPI$ with
the additional requirement that $P \wbPI Q$ implies 
\begin{itemize}
\item
$P \Uparrow $ iff  $Q \Uparrow $ \hfill $(*)$
\end{itemize} 
where the predicate $ \Uparrow $ holds for an  agent $R$  if it
has a divergence (an infinite path 
consisting of only $\tau$ actions). 

\begin{lemma}
\label{l:bisDIV_op}
Relation $\wbPIDIV$ implies $\wbOPI$.
\end{lemma} 
\noindent
The same result applies to  other relations; e.g., the `expansion  respecting divergence'
(defined by adding the clause $(*)$ above to those of the expansion
relation $\expa$)
 on the ordinary LTS implies expansion  on the \opLTS.

A comment about expansion ($\expa$) is worthwhile. 
As pointed out earlier,   the use of expansion
is important in proofs  about operational correspondence results. 
Such uses are  due to 
the appearance of expansion in 
 lemmas such as  
 Lemma~\ref{l:opt_sound} (and Lemmas~\ref{l:opt_stuck}-\ref{l:opt_con}, in
 Appendix~\ref{a:oc}),
 and
come from application of some
simple algebraic laws that respect  divergence (i.e., the laws do not add or remove any
divergence). Indeed the laws either are valid for strong bisimilarity, or they are laws
such as 
\[ 
{\res{a} ( \inp a {\tilx}.P | \bout a \tilx . Q)} \;  \contr \;
{\res{a} (P |  Q)  }
 \]
As a consequence, all occurrences of $\expa$ 
%in such lemmas 
%(notably in Lemmas~\ref{l:opt_sound}-\ref{l:opt_con})
can be replaced by the expansion relation defined on the \opLTS ($\expaOP$). 
% This in turn implies that the same can be done   for 
%  all occurrences of $\expa$ in
% Theorems~\ref{t:opcorrCON_w} and \ref{t:opcorrCON2_w}, and in Lemma~\ref{l:opcorrCON3_w}
% (as they rely on Lemmas~\ref{l:opt_sound}, 
% and Lemmas~\ref{l:opt_stuck}-\ref{l:opt_con}).
In this manner, for instance, we can   refine (the proof of)  Lemma~\ref{l:opcorrCON3_w}
and derive: 

 \begin{theorem} %[from \piI to  OGS, sequential LTS]
 \label{t:opcorrCON3_w}
$ $ \begin{enumerate}

% \item  If
% $\encoConV{\FF } \arrN  {\initName (p)} \PP$, then there is $\FF'$ 
% such that
%  $\FF \arr{\initName (p)} \FF'$ 
% and $\PP = \encoConV{\FF' }$ ;

\item If 
$ \encoConI \FF \ArrN {} \PP$ then there is $\FF'$ such that 
$\FF \LongrightarrowA \FF'$  and  
$\PP  \wbOPI 
 \encoConI {\FF'}$ ; 


\item If
$ \encoConI \FF \ArrN\ell \PP$
 then there is $\FF'$ such that
 $\FF \ArrA\ell \FF'$
and $\PP \wbOPI \encoConI {\FF'}$.
\end{enumerate} 
\end{theorem} 

\noindent
A final remark  about 
behavioural equivalences on the ordinary and \opLTS
concerns the encoding of $\lambda$-terms. 
On such terms, 
the ordinary
bisimilarity $\wbPI$ implies the `bisimilarity respecting divergence' $\wbPIDIV$. This
is because, intuitively, a term  $\encoIa {M} p$ is divergent iff $M$ is so in the call-by-value
$\lambda$-calculus. 
(Again, this property boils down to the fact that the uses of
$\expa$ in Lemmas such as~\ref{l:opt_sound}-\ref{l:opt_con} 
 respect divergence.)
   See Appendix~\ref{a:behav} for   details. 


From Theorem~\ref{t:opcorrCON3_w} we derive: 

\DSOCTe



\begin{corollary}
\label{c:traces}
For any configuration $\FF$ and trace $s$, we have 
  $\FF \ArrA s$   iff 
 $ \encoEnvI \FF  \ArrN s$.
\end{corollary} 

\begin{remark}
\label{r:bn2} 
% {\em 
We recall that, following Remark~\ref{r:bn} on the usage of bound names, 
in Corollary~\ref{c:traces} 
the bound names in $s$ are fresh;   thus they do not appear in $\FF$.
(Similarly in Theorems~\ref{t:opcorrCON2_w} and 
\ref{t:opcorrCON3_w}, for 
 the bound names
in $\ell$.)
% }
\end{remark} 


 \begin{remark}
 \label{r:expa_op_ord} 
 Corollary~\ref{c:traces} relies on Theorems~\ref{t:opcorrCON_w} and \ref{t:opcorrCON3_w}.
 The corollary talks about the \opLTS of \piI; however the theorems make use of the
 ordinary expansion relation  $\expa $, which is defined on the ordinary LTS.  Such uses of
 expansion can however be replaced by expansion on the \opLTS (defined as ordinary
 expansion, but on the \opLTS). For more details on this, see
 Appendix~\ref{a:behav}. 
 \end{remark} 

\noindent
As a consequence of Corollary~\ref{c:traces}, trace equivalence is the same, 
on \AOGS configurations and on the encoding of \piI terms. 
Moreover, from Theorem~\ref{t:opcorrCON3_w} the same result holds under a bisimulation semantics.
Further,
since the LTS produced by \AOGS is deterministic, 
its trace semantics coincides with its bisimulation semantics.
We can thus conclude as in Corollary~\ref{c:fa_traces_bis}.  
We recall that $\TEN$ and $\wbOPI$ are, respectively, trace equivalence and bisimilarity 
 between \piI processes in the \opLTS; similarly for  
$\TEA$ and $\wbA$ between \AOGS terms.  

\begin{corollary}
%[full abstraction for  traces and bisimilarity, between Alternating OGS and \piI with \opLTS] %,  \opLTS]
\label{c:fa_traces_bis}
%$ $ 
For any $\FF,\FF'$ we have:
% \begin{itemize}
% \item
 $\FF\TEA \FF'$ iff 
$ \encoConI \FF \TEN \encoConI {\FF'}$
iff
 % \item
 $\FF\wbA \FF'$ iff 
$ \encoConI \FF \wbOPI \encoConI {\FF'}$.
%\end{itemize}
 \end{corollary} 

\noindent
Corollary~\ref{c:fa_traces_bis}  holds in particular when $\FF$ is the initial configuration
for a $\lambda$-term.  That is, 
the equality induced on call-by-value $\lambda$-terms by  
their representation in \AOGS and in \piI 
(under the \opLTS) is the same, both employing traces and employing bisimulation 
to handle   the observables for the two
models. 
\begin{corollary}
\label{c:OGS_opLTSpiI}
For any $\lambda$-terms $M,N$,  we have: 
\\
$\conf{M} \TEA \conf {N}$ iff
$\conf{M} \wbA \conf {N}$ iff
$\encoEnvI {{M} } \TEN\encoEnvI {{N}}$ iff 
$\encoEnvI {{M}} \wbOPI \encoEnvI {{N}}$. 
\end{corollary} 

\noindent
From Theorem~\ref{t:opcorrCON2_w} and Corollary~\ref{c:traces}, it also follows that  $\FF$ and
$\encoConI \FF$ are weakly bisimilar,  on the union of the respective  LTSs.
